\section{The quartic spectral flow}\label{inf:sec-flow}

Throughout this section $p\ge64$, $\varsigma\in I=[1/64,1/16]$,
$r=1/2-\sqrt \varsigma$, $R\in[2p-1,2p+4]$, and $0<\tau=|\theta|\le20$.
Unless indicated otherwise, constants are absolute and independent of
$p,\varsigma,R,\theta$. The fixed window used for the flow argument is
$W_{\mathrm{flow}}$. The coefficient inverse uses the separately
defined windows in Appendix~\ref{inf:sec-notation}.

\subsection{Angular matrices and their weighted variation}

Identify the angular Hilbert spaces for different $\varsigma$ by their
orthonormal bases $P_j(\varsigma)$. In these coordinates the angular operator
is the fixed diagonal matrix $\Lambda=\diag(\lambda_j)$, and
multiplication by $x$ is the matrix $X_\varsigma$ in
\eqref{inf:jacobi-matrix}. Define
\[
 T_\varsigma=\frac{\mathsf H_2(X_\varsigma)}{pg_2},\qquad
 w_j=\frac{j+1}{p+j+1},\qquad \mathsf W=\diag(w_j).
\]
The same notation $T_\varsigma(x)$ denotes its scalar polynomial before
identification. All estimates involving $\mathsf W^{-1}$ initially
mean estimates on finite sequences, extended by boundedness.

\begin{lemma}[Weighted angular bounds]\label{inf:angular-weighted}
For $k=0,1,2$, the matrix $\partial_\varsigma^kT_\varsigma$ has bandwidth two and
\begin{equation}\label{inf:T-entry-bound}
 |(\partial_\varsigma^kT_\varsigma)_{ij}|\le C\min(w_i,w_j),\qquad
 \nrm{(\partial_\varsigma^kT_\varsigma)\mathsf W^{-1}}+
 \nrm{\mathsf W^{-1}(\partial_\varsigma^kT_\varsigma)}\le C.
\end{equation}
The matrices $X_\varsigma$ and their first two derivatives are bounded uniformly.
If
\begin{equation}\label{inf:G-angular}
 G_\varsigma=4x(1-x)(\partial_xT_\varsigma)^2
     =\tfrac12[T_\varsigma,[\Lambda,T_\varsigma]],
\end{equation}
then, for $k=0,1,2$,
\begin{equation}\label{inf:G-weighted}
 \nrm{\mathsf W^{-1}(\partial_\varsigma^kG_\varsigma)\mathsf W^{-1}}\le Cp.
\end{equation}
\end{lemma}
\begin{proof}
From \eqref{inf:recurrence-a}, uniformly for $j\ge1$ in the stated range,
\[
 \mathfrak a_j^2\le C\frac j{p+j},\qquad
 \frac{\partial_\varsigma\mathfrak a_j^2}{\mathfrak a_j^2}
 =-\frac{p^2}{(j+p/2-1)^2-p^2\varsigma}.
\]
For $j\ge1$, write
$\mathfrak D_j=(j+p/2-1)^2-p^2\varsigma$. The exact derivatives are
\begin{equation}\label{inf:angular-second-derivative}
 \frac{\partial_\varsigma\mathfrak a_j}{\mathfrak a_j}
 =-\frac{p^2}{2\mathfrak D_j},\qquad
 \frac{\partial_\varsigma^2\mathfrak a_j}{\mathfrak a_j}
 =-\frac{p^4}{4\mathfrak D_j^2}.
\end{equation}
The factor $\mathfrak D_j$ is bounded below by $c(p+j)^2$,
so $|\partial_\varsigma^k\mathfrak a_j|\le C\mathfrak a_j$
for $k=1,2$. The coefficient $\mathfrak a_0$ and both its
derivatives are identically zero. The diagonal entries $\mathfrak b_j$ have uniformly
bounded first two derivatives: $\varsigma$ stays away from zero, so
$\partial_\varsigma^k\sqrt \varsigma$ is bounded for $k\le2$.
The two-step entry of $\mathsf H_2(X_\varsigma)$ is
$\mathfrak a_{j+1}\mathfrak a_{j+2}$ and has size, together with its
two derivatives, at most $Cw_j$.

The one-step entry is
\[
 \mathfrak a_{j+1}\left(\mathfrak b_j+\mathfrak b_{j+1}
                               -\frac{2(\ab+1)}{p+2}\right).
\]
Its parenthesis equals $(r-1/2)$ times
\[
 p(p-2)\left\{
 \frac1{(p+2j-2)(p+2j)}+
 \frac1{(p+2j)(p+2j+2)}\right\}-\frac{2p}{p+2}.
\]
The braces with their prefactor have exactly the compensating value
at $j=0$. For $j\le p$, subtract this value and differentiate the
reciprocal products in $j$ to bound the difference by $Cj/p$.
For $j\ge p$ each term is bounded by a constant. Thus the
parenthesis and both its $\varsigma$ derivatives are $O(j/(p+j))$.
Multiplication by $\mathfrak a_{j+1}$ preserves the bound $Cw_j$.

The diagonal entries of $\mathsf H_2(X_\varsigma)$ are
$M_j=c_j^{\mathrm{diag}}+d_j^{\mathrm{diag}}\varsigma$, where
\begin{align}
 c_j^{\mathrm{diag}}&=\frac{j(p-1)(j+p-1)}{2(p+1)(2j+p-3)(2j+p+1)},
 \label{inf:T-diagonal}\\
 d_j^{\mathrm{diag}}&=-\frac{2jp^2(j+p-1)(p^2+7p-8j^2-8jp+8j)}
 {(2j+p)(p+1)(p+2)(2j+p-3)(2j+p-2)(2j+p+1)}.\nonumber
\end{align}
These formulas follow by the multiplication calculation in the proof
of Lemma~\ref{inf:quartic-normalization}, now at general $j$.
The denominator factors are comparable to $p$ or $p+j$;
hence $|c_j^{\mathrm{diag}}|+|d_j^{\mathrm{diag}}|\le Cj/(p+j)$.
The scalar $pg_2$ is bounded away from zero and has bounded derivatives
of every fixed order. Division by it proves the entry bounds.
For $|i-j|\le2$ the weights have ratios between $1/3$ and $3$.
The row and column sums over the five diagonals prove
\eqref{inf:T-entry-bound}. The same row-sum argument for $X_\varsigma$
and its derivatives proves their asserted bounds.

The identity in \eqref{inf:G-angular} follows by expanding
$[\Lambda,T]= (\Lambda T)_{\mathrm{mult}}
-8x(1-x)T_x\partial_x$.
The function $G_\varsigma$ is a degree-four polynomial in $X_\varsigma$, with
coefficients and their first two derivatives uniformly bounded.
Its bandwidth is four and its unweighted norm is $O(1)$.
When both indices exceed $p-4$, the weights are bounded below and
this already proves the required bound.
For the other nonzero entries all indices are at most $p+4$, and
\[
 (G_\varsigma)_{ij}=\frac12\sum_l(2\lambda_l-\lambda_i-\lambda_j)
                                 (T_\varsigma)_{il}(T_\varsigma)_{lj}.
\]
Only finitely many neighbours occur; their eigenvalue differences
are $O(p)$ and each product is at most $Cw_iw_j$ by
\eqref{inf:T-entry-bound}. The product rule gives the same estimate
for two derivatives. Summing over the nine diagonals proves
\eqref{inf:G-weighted}.
\end{proof}

\begin{lemma}[Positive angular variation]\label{inf:angular-positive}
For every subset $\mathcal I\subset\{8,9,\ldots\}$ containing no
adjacent indices, let $T_{\mathcal I}$ denote the principal angular
submatrix with those indices. Then
\begin{equation}\label{inf:angular-positive-bound}
 T'_{\mathcal I}(\varsigma)\ge c_0\mathsf W_{\mathcal I}^2,
 \qquad c_0=10^{-4},
\end{equation}
for $p\ge64$. The inequality holds for
finite or infinite $\mathcal I$ as a bounded-operator inequality.
\end{lemma}
\begin{proof}
Differentiate
$M_j/g_2=(c_j^{\mathrm{diag}}+d_j^{\mathrm{diag}}\varsigma)/
(c_2^{\mathrm{diag}}+d_2^{\mathrm{diag}}\varsigma)$.
The numerator is
\begin{align}
 \mathcal N_j&=d_j^{\mathrm{diag}}c_2^{\mathrm{diag}}
                    -c_j^{\mathrm{diag}}d_2^{\mathrm{diag}}\nonumber\\
 &=\frac{24jp^3(j-2)(p-1)(j+p-1)(j+p+1)}
 {(2j+p)(p+1)(p+2)^2(p+4)(p+5)
  (2j+p-3)(2j+p-2)(2j+p+1)}.
 \label{inf:N-angular}
\end{align}
Define also $\partial_\varsigma(\mathfrak a_j^2/g_2)
=\mathcal E_j/g_2^2$, where
\begin{align}
 \mathcal E_j&=\frac{4jp^2(j+p-2)F(p,j)}
 {(p+1)(p+2)^2(p+4)(p+5)(2j+p-3)(2j+p-2)^2(2j+p-1)},
 \label{inf:E-angular}\\
 F(p,j)&=(p^2-9p-16)j^2+(p^3-11p^2+2p+32)j
                         -5p^3+3p^2+8p-12.\nonumber
\end{align}
These identities are obtained by substituting
\eqref{inf:T-diagonal}, \eqref{inf:recurrence-a}, and
\eqref{inf:g2-exact} and clearing their positive denominators.
For $u=p-64$, $j_{\mathrm{pos}}=j-7$ one has the explicit identity
\begin{align*}
 F(p,j)={}&(u^2+119u+3504)j_{\mathrm{pos}}^2\\
 &+(u^3+195u^2+12548u+266304)j_{\mathrm{pos}}\\
 &+2u^3+359u^2+20957u+394500.
\end{align*}
Thus $\mathcal E_j>0$ for $p\ge64,j\ge7$.
The polynomial identity in Appendix~\ref{inf:sec-positivity} gives
\begin{equation}\label{inf:half-diagonal}
 \frac{\mathcal N_j}{2}
 >\frac23(\mathcal E_{j-1}+\mathcal E_j+
                 \mathcal E_{j+1}+\mathcal E_{j+2}),\qquad j\ge8.
\end{equation}
All denominators there are positive on the stated range.

The adjacent recurrence coefficients satisfy
$3/4<\mathfrak a_{j+1}/\mathfrak a_j<4/3$ for $j\ge7$.
For the upper bound discard the denominator ratio in
\eqref{inf:recurrence-a} to obtain
\[
 \frac{\mathfrak a_{j+1}^2}{\mathfrak a_j^2}
 \le\frac87\left(\frac{23}{22}\right)^2\frac{70}{69}<\frac{16}{9}.
\]
For the inverse ratio discard the increasing numerator factors and
use $2j+p\ge78$ to get
\[
 \frac{\mathfrak a_j^2}{\mathfrak a_{j+1}^2}
 \le\frac{78^2\cdot79}{76^2\cdot75}<\frac{16}{9}.
\]
The corresponding rational bound decreases with $2j+p$, factor by
factor, so these estimates cover every larger $j,p$.
On one parity the only off-diagonal entries of $T_\varsigma$ are
$\mathfrak a_{j+1}\mathfrak a_{j+2}/(pg_2)$.
Their derivatives are positive and, for the upper neighbour, at most
\[
 \frac{2}{3pg_2^2}(\mathcal E_{j+1}+\mathcal E_{j+2}).
\]
Indeed differentiate the product as one half the sum of the two
square derivatives times the adjacent coefficient ratios.
Equation~\eqref{inf:half-diagonal} and
$2|v_jv_k|\le|v_j|^2+|v_k|^2$ now give
\[
 \langle T'_\varsigma v,v\rangle\ge
 \sum_j\frac{\mathcal N_j}{2pg_2^2}|v_j|^2
\]
on either parity and every subset thereof.
For $8\le j\le p$, the numerator and denominator in
\eqref{inf:N-angular} are respectively at least $9j^2p^6$ and at
most $3456p^9$. Since $g_2<1/p$, the margin is at least
$j^2/(768p^2)$. For $j\ge p$, the corresponding estimates
$6p^4j^4$ and $3456p^5j^4$ give a margin at least $1/1152$.
Together with $w_j\le(9/8)j/p$ for $j\ge8$ and $w_j\le1$, these
bounds imply \eqref{inf:angular-positive-bound} with $c_0=10^{-4}$.

A set with no adjacent indices has no opposite-parity off-diagonal
entry: the matrix has bandwidth two. Its principal submatrix is
therefore the direct sum of the two parity submatrices.
The bounds extend from finite vectors by density and
Lemma~\ref{inf:angular-weighted}.
\end{proof}

\subsection{The common Dirichlet form}

Let $\zeta\in(0,1)$ be nonsquared reference radius and put
\[
 \epsilon_{\mathrm q}=-\frac{\theta}{2R^2},\qquad
 \rho_\varsigma(x)=1+\epsilon_{\mathrm q}T_\varsigma(x),\qquad c_p=p-\tfrac12.
\]
Use the Hilbert space
$\Hc_{\mathrm{flow}}=L^2((0,1),d\zeta)\otimes\ell^2(\mathbb N_0)$,
with the Jacobi identification already made. The ball operator and
its positive shift are
\begin{equation}\label{inf:flow-ball-operators}
 H_0=-\partial_\zeta^2+
          \frac{c_p(c_p-1)+\Lambda}{\zeta^2},\qquad
 K_0=H_0+R^2,\qquad
 \mathsf D=\zeta\partial_\zeta+\tfrac12.
\end{equation}
This is the mass-unitary representation of the operator already
denoted by $H_0$. Its spectrum and all its spectral projections
are unchanged by the unitary map.

\begin{proposition}[Exact pullback form]\label{inf:flow-form}
For the weighted quartic domain $0<\varrho<\rho_\varsigma(x)$, let
$U(\zeta,x)=u(\zeta\rho_\varsigma(x),x)$ and
$v=\sqrt{2p}\,\zeta^{p-1/2}\rho_\varsigma^pU$. Put
$L_\varsigma=\partial_x\log\rho_\varsigma$, $\alpha_\varsigma=\rho_\varsigma^{-2}$, and
$q_x=4x(1-x)$. Its Dirichlet form is
\begin{equation}\label{inf:flow-formula}
 \mathfrak q_\varsigma[v,w]=\int\alpha_\varsigma\left\{
 (v_\zeta-c_pv/\zeta)\overline{(w_\zeta-c_pw/\zeta)}
 +\frac{q_x}{\zeta^2}(v_x-L_\varsigma\mathsf Dv)
                       \overline{(w_x-L_\varsigma\mathsf Dw)}\right\}
 \,d\zeta\,d\mu_r.
\end{equation}
Define, initially as weak expressions on the common core,
\begin{align}
 \mathcal L(M)&=MH_0+\frac1{2\zeta^2}[\Lambda,M]\mathsf D,
 \label{inf:linear-form-operator}\\
 \mathfrak g_\gamma[v,w]&=\int\zeta^{-2}\gamma
                                  (\mathsf Dv)\overline{\mathsf Dw}.
 \nonumber
\end{align}
Write $\mathcal G(\gamma)$ for the weak form operator characterized
on this core by
$\langle\mathcal G(\gamma)v,w\rangle=\mathfrak g_\gamma[v,w]$.
The operator associated with $\mathfrak q_\varsigma$ is, in the sense of forms,
\begin{equation}\label{inf:flow-exact-operator}
 H_\varsigma=H_0+\mathcal L(\alpha_\varsigma-I)+\mathcal G(\gamma_\varsigma),
 \qquad
 \gamma_\varsigma=\epsilon_{\mathrm q}^2
                   (I+\epsilon_{\mathrm q}T_\varsigma)^{-4}G_\varsigma.
\end{equation}
The exact form has the same closed domain as $H_0$ for large $p$.
The family is $C^1$ in the relative form norm and has compact
resolvent. At integer block sizes it is unitarily equivalent to the
invariant Euclidean Dirichlet Laplacian on $\Omega_{\mathrm q}/R$.
\end{proposition}
\begin{proof}
With the fixed normalization by the unit-ball volume, the radial
change of variables has density $2p\zeta^{2p-1}\rho_\varsigma^{2p}$,
proving the mass normalization.
Its angular derivative is $U_x-\zeta L_{\varsigma}U_\zeta$.
Substitution gives
\[
 U_x-\zeta L_{\varsigma}U_\zeta
 = (2p)^{-1/2}\zeta^{-c_p}\rho_\varsigma^{-p}
       \{v_x-L_\varsigma(\zeta v_\zeta+v/2)\}.
\]
The radial derivative gives
$(2p)^{-1/2}\zeta^{-c_p}\rho_\varsigma^{-p}(v_\zeta-c_pv/\zeta)$.
These identities prove \eqref{inf:flow-formula}.

To integrate the form, radial integration of its first square gives
$\alpha_\varsigma(-\partial_\zeta^2+c_p(c_p-1)/\zeta^2)$.
The purely angular square gives
\[
 \zeta^{-2}\{\alpha_\varsigma\Lambda+\tfrac12[\Lambda,\alpha_\varsigma]
                         -\tfrac12(\Lambda\alpha_\varsigma)_{\mathrm{mult}}\}.
\]
For the cross terms use $\alpha_{\varsigma}L_\varsigma=-(\alpha_\varsigma)_x/2$ and
$\mathsf D^*=-\mathsf D+2$ in radial measure $\zeta^{-2}d\zeta$.
Their sum is
\[
 \zeta^{-2}\{\tfrac12[\Lambda,\alpha_\varsigma](\mathsf D-1)
                         +\tfrac12(\Lambda\alpha_\varsigma)_{\mathrm{mult}}\}.
\]
The multiplication terms cancel. The remaining square has
coefficient $\alpha_{\varsigma}q_xL_\varsigma^2$, giving
\eqref{inf:flow-exact-operator} and \eqref{inf:G-angular}.
These are identities on finite angular sums with radial functions
compactly supported in $(0,1)$.

We specify the closure. Let $\mathcal V_0=\Dom(K_0^{1/2})$, defined
as the completion of that core in its positive form norm.
Angular spectral truncations commute with $K_0$ and converge in
this norm. In each remaining radial block the Friedrichs definition
gives approximation by compactly supported radial functions.
Thus the stated core is form dense. Its angular integration by
parts has no endpoint term: the boundary factor is
$x^{\ab}(1-x)^{\bb}$ and the test functions are polynomial there.
Cutting off a smooth angular function near $x=0$ costs
$O(\delta^{\ab-1})$ in energy, and near $x=1$ costs
$O(\delta^{\bb-1})$. Both tend to zero since $\ab,\bb>1$.
At the physical origin a radial cutoff costs $O(\delta^{2p-2})$.
Consequently the block axes and origin have zero capacity for
this weighted energy.
For integer blocks this is the invariant Euclidean $H^1_0$
realization, by averaging smooth Dirichlet approximations over the
compact symmetry group. For real parameters define the original weighted Dirichlet
realization as the closure of smooth functions compactly supported
in the open meridian rectangle, using the norm from mass plus
\[
 2p\int_0^1\!\int_0^1
 \left(|U_\zeta|^2+\zeta^{-2}q_x|U_x|^2\right)
                  \zeta^{2p-1}\,d\zeta\,d\mu_r.
\]
The polynomial angular core has the same closure. The cutoffs
just estimated put every finite angular
polynomial times a compactly supported smooth radial function
in this closure. Conversely, for a smooth function compactly
supported in the meridian rectangle, let $U_J$ be its angular
projection on $P_0,\ldots,P_J$.
Integration by parts in $x$ gives
$\langle\Lambda_r U,P_j\rangle=\lambda_j\langle U,P_j\rangle$.
The smooth compact support makes $\Lambda_rU$ square integrable,
so Parseval gives square summability of these coefficients with
weight $\lambda_j^2$. It follows that $U_J\to U$ in angular
energy as well as in angular $L^2$.
The same projection commutes with the radial derivative;
Parseval applied to $U_\zeta$ gives convergence of its radial
energy. The radial support stays in a compact subinterval of
$(0,1)$, so the factor $\zeta^{-2}$ is bounded there.
Dominated convergence in the radial integral proves convergence
in the full weighted form norm. Each projected radial coefficient
is smooth with compact support. Thus both cores have the same
closure, establishing the real-parameter domain.

For integer blocks, the identification with the Euclidean
invariant form domain can also be read in both directions.
Average a compactly supported smooth Euclidean approximation
over $O(a)\times O(b)$; averaging is a contraction for mass and
Dirichlet energy. The result is smooth and invariant.
Cut it off at the origin and at each block axis. In squared
angular distance the extra energy is bounded by
$C\delta^{\ab-1}+C\delta^{\bb-1}$; at the origin it is bounded
by $C\delta^{2p-2}$. All tend to zero. Away from those sets it
is a smooth function of the meridian coordinates with compact
support, so it belongs to the weighted closure just described.
In the reverse direction a meridian test function with compact
support away from the axes and origin lifts to a smooth
invariant Euclidean test function, extended by zero across
those sets. This proves equality of the closed domains.

The estimates proved immediately below show that the difference of
\eqref{inf:flow-formula} and the ball form has relative norm below
$1/2$ on $\mathcal V_0$. Hence its closed form norm is equivalent to
that of $K_0$, and its closure has exactly domain $\mathcal V_0$.
The angular mass factor $\rho_\varsigma^p$ also preserves the
weighted first-order domain. The required control at the origin
is Hardy's inequality in the radial weight:
\begin{equation}\label{inf:flow-Hardy}
 \int_0^1|U|^2\zeta^{2p-3}\,d\zeta
 \le\frac1{(p-1)^2}
      \int_0^1|U_\zeta|^2\zeta^{2p-1}\,d\zeta.
\end{equation}
For compactly supported $U$ this follows by integrating
$(\zeta^{2p-2})'|U|^2$ and applying Cauchy--Schwarz; completion
extends it. Integrate also against $\mu_r$.
Multiplication by $\rho_\varsigma^{\pm p}$ differentiates in $x$
by the extra term $\pm p L_\varsigma U$. Its angular energy has
the factor $p^2|U|^2/\zeta^2$, which is controlled by
\eqref{inf:flow-Hardy}. Both multipliers and their inverses are
bounded for fixed parameters. Thus the unitary mass transformation
preserves the closure, including its origin behaviour.
The radial graph map is bi-Lipschitz even at the origin: on the
sphere its radial factor and first derivatives are bounded and it
is bounded away from zero. Such a map transports Dirichlet energy
and its closure. Therefore the just described closure is the
physical one, including at the origin. The relative derivative
bounds below prove norm $C^1$ dependence. Compactness of the
embedding $\mathcal V_0\hookrightarrow\Hc_{\mathrm{flow}}$ persists
under equivalent form norms and proves compact resolvent.
\end{proof}

\begin{lemma}[Relative form estimates]\label{inf:flow-relative}
Set $B_{1,\varsigma}=-2\mathcal L(\epsilon_{\mathrm q}T_\varsigma)$ and
$\mathcal R_{2,\varsigma}=H_\varsigma-H_0-B_{1,\varsigma}$. For $k=0,1,2$ and large $p$,
\begin{align}
 \nrm{K_0^{-1/2}(\partial_\varsigma^kB_{1,\varsigma})K_0^{-1/2}\mathsf W^{-1}}
 &\le C\tau p^{-2},\label{inf:flow-linear-relative}\\
 \nrm{\mathsf W^{-1}K_0^{-1/2}(\partial_\varsigma^k\mathcal R_{2,\varsigma})
               K_0^{-1/2}\mathsf W^{-1}}
 &\le C\theta^2p^{-3}.\label{inf:flow-quadratic-relative}
\end{align}
These are estimates on the full radial spaces. They also establish
the closure and differentiability assertions in
Proposition~\ref{inf:flow-form}.
\end{lemma}
\begin{proof}
We first prove the band estimate used to sum the exact form.
If an angular matrix $M$ has bandwidth $m$ and
$|M_{ij}|\le M_*w_i^aw_j^b$, $a,b\in\{0,1\}$, then
\begin{equation}\label{inf:band-form-bound}
 \nrm{\mathsf W^{-a}K_0^{-1/2}\mathcal L(M)
                K_0^{-1/2}\mathsf W^{-b}}
 \le C(m+1)^3M_*.
\end{equation}
Indeed the radial forms
$K_{0,j}=-\partial_\zeta^2+(\nu_j^2-1/4)/\zeta^2+R^2$
are comparable with ratio at most $C(m+1)^2$ for $|i-j|\le m$.
This follows from $\nu_i/\nu_j$ and its inverse being at most
$C(m+1)$. Thus
$\nrm{K_{0,i}^{-1/2}K_{0,j}^{1/2}}\le C(m+1)$.
The block $M_{ij}H_{0,j}$ in \eqref{inf:linear-form-operator}
is bounded after normalization since
$\nrm{H_{0,j}K_{0,j}^{-1}}\le1$.
For the commutator block use the two form bounds
\[
 \nrm{\zeta^{-1}K_{0,i}^{-1/2}}\le C/(p+i),\qquad
 \nrm{\zeta^{-1}\mathsf DK_{0,j}^{-1/2}}\le C.
\]
The first follows from the centrifugal potential; the second from
$\zeta^{-1}\mathsf D=\partial_\zeta+1/(2\zeta)$ and the same
positive form. Also
$|\lambda_i-\lambda_j|\le Cm(p+i+j)\le Cm(m+1)(p+i)$.
Each normalized block is bounded by $C(m+1)^2M_*$ after its
weights are removed. The row and column sums involve at most
$2m+1$ blocks. The scalar Schur inequality proves
\eqref{inf:band-form-bound} on the form core, then by closure.

For $n\ge2$, Lemma~\ref{inf:angular-weighted} gives
\[
 |(\partial_\varsigma^kT_\varsigma^n)_{ij}|\le n^2C^nw_iw_j,\qquad k\le2,
\]
with the factor $n^2$ replaced by one when $k=0$ if desired.
Put one inverse weight on the first factor and one on the last;
each differentiated endpoint has the same one-sided bound, and
interior factors are uniformly bounded. The product rule has at
most a fixed multiple of $n^2$ terms. The bandwidth is $2n$.
The series
\[
 \alpha_\varsigma=I-2\epsilon_{\mathrm q}T_\varsigma
       +\sum_{n\ge2}(-1)^n(n+1)\epsilon_{\mathrm q}^nT_\varsigma^n
\]
therefore has a two-sided weighted remainder under
\eqref{inf:band-form-bound} of size $C\epsilon_{\mathrm q}^2$,
including two derivatives. Its dominating series is
$\sum n^6(C|\epsilon_{\mathrm q}|)^n$, which converges uniformly.

Equation~\eqref{inf:G-weighted} and the convergent series for
$(I+\epsilon_{\mathrm q}T_\varsigma)^{-4}$ give
\[
 \nrm{\mathsf W^{-1}(\partial_\varsigma^k\gamma_\varsigma)\mathsf W^{-1}}
 \le Cp\epsilon_{\mathrm q}^2,\qquad k\le2.
\]
To move the left weight through the inverse series, use boundedness
of $\mathsf W^{-1}T_\varsigma\mathsf W$ and its derivatives, which follows
from \eqref{inf:T-entry-bound}. Since
$\zeta^{-1}\mathsf D$ is controlled by $K_0^{1/2}$ and commutes
with the angular weights, the same bound holds for the normalized
form $\mathcal G(\gamma_\varsigma)$.
Now $|\epsilon_{\mathrm q}|\le C\tau/p^2$.
The first-order band estimate proves
\eqref{inf:flow-linear-relative}; the two remainder estimates prove
\eqref{inf:flow-quadratic-relative}.
Their unweighted sum is below $1/2$ for large $p$.

The $k=2$ bounds provide a uniform modulus of continuity for the
first derivative: Taylor's integral formula on the finite bands,
followed by the uniformly convergent series, bounds the difference
quotient remainder in relative form norm. Thus the family is norm
$C^1$ on the common form space. These bounds also prove the
closure and regularity assertions of Proposition~\ref{inf:flow-form}.
\end{proof}

\subsection{The boundary matrix and the transport term}

Fix a finite $W_{\mathrm{flow}}$ and set
$\PP=\one_{[-W_{\mathrm{flow}},W_{\mathrm{flow}}]}(H_0-R^2)$.
For large $p$, Lemma~\ref{inf:window-isolation} says its states have
distinct, nonadjacent angular indices. Index its orthonormal states by
$i$, let $j(i)$ be their angular indices and $z_i$ their positive
zeros, and write
\[
 D_\PP=\diag(z_i^2-R^2),\quad
 Z_\PP=\diag(z_i/R),\quad
 T_\PP^{\mathrm{ang}}=[(T_\varsigma)_{j(i),j(i')}],\quad
 \mathcal S_\varsigma=\PP(\epsilon_{\mathrm q}T_\varsigma\mathsf D)\PP.
\]
The matrix $T_\PP^{\mathrm{ang}}$ is the \emph{angular principal
submatrix}, while $\mathcal S_\varsigma$ is the actual Hilbert-space
compression with radial overlaps.

\begin{lemma}[Transport identity]\label{inf:transport}
One has $\mathcal S_\varsigma^*=-\mathcal S_\varsigma$ and the exact identity
\begin{equation}\label{inf:transport-identity}
 \PP B_{1,\varsigma}\PP=\theta Z_\PP T_\PP^{\mathrm{ang}} Z_\PP
                             -[D_\PP,\mathcal S_\varsigma].
\end{equation}
For $k=0,1$,
\begin{equation}\label{inf:transport-bounds}
 \nrm{\PP(\partial_\varsigma^kB_{1,\varsigma})\PP\mathsf W^{-1}}\le C\tau,
 \qquad
 \nrm{(\partial_\varsigma^k\mathcal S_\varsigma)\mathsf W^{-1}}\le C\tau/p,
\end{equation}
with constants independent of $W_{\mathrm{flow}}$ once
$p^2\ge W_{\mathrm{flow}}$.
\end{lemma}
\begin{proof}
As interior differential expressions,
$[H_0,\mathsf D]=2H_0$ and
$[H_0,h]=\zeta^{-2}[\Lambda,h]$ for an angular multiplier $h$.
Thus $B_{1,\varsigma}=-[H_0,\epsilon_{\mathrm q}T_\varsigma\mathsf D]$.
Green's identity applied to two ball Dirichlet eigenfunctions gives
\[
 \langle\phi_i,H_0(h\mathsf D\phi_{i'})\rangle
 =z_i^2\langle\phi_i,h\mathsf D\phi_{i'}\rangle
 +\int h\,\partial_\zeta\phi_i(1)
                 \partial_\zeta\phi_{i'}(1)\,d\mu_r.
\]
The origin term vanishes by the regular radial behaviour. The
boundary derivatives are $\sqrt2z_iP_{j(i)}$ by
Lemma~\ref{inf:bessel-spacing}. Substitution of
$h=-\theta T_\varsigma/(2R^2)$ gives
\eqref{inf:transport-identity}. Radial integration on pairs of
Dirichlet eigenfunctions gives $\mathsf D^*=-\mathsf D$; the
angular multiplier commutes with $\mathsf D$. Hence $\mathcal S_\varsigma$ is
skew-adjoint.
On the fixed window $\nrm{K_0^{1/2}\PP}\le Cp$.
Equation~\eqref{inf:flow-linear-relative} proves the first bound
in \eqref{inf:transport-bounds}. The form estimate also gives
$\nrm{\mathsf D\PP}\le Cp$; use
\eqref{inf:T-entry-bound} and
$|\epsilon_{\mathrm q}|\le C\tau/p^2$ for the second bound.
\end{proof}

\subsection{Eliminating the full complementary space}

\begin{lemma}[Complementary inverse and remainder]
\label{inf:flow-complement}
There is an absolute lower bound on $W_{\mathrm{flow}}\ge100$ such
that the following holds for $p$ sufficiently large depending on
that fixed window. For every $\varsigma\in I$ and real $|\lambda|\le1$,
the form operator
$\PP^\perp(H_\varsigma-R^2-\lambda)\PP^\perp$ is invertible. Its exact
Schur matrix is
\begin{equation}\label{inf:flow-schur}
 F_\PP(\varsigma,\lambda)=D_\PP-\lambda+
                        \PP B_{1,\varsigma}\PP+\mathcal R_\varsigma(\lambda),
\end{equation}
where, for $k=0,1$,
\begin{equation}\label{inf:flow-schur-remainder}
 \nrm{\mathsf W^{-1}(\partial_\varsigma^k\mathcal R_\varsigma)\mathsf W^{-1}}
 \le C\theta^2(W_{\mathrm{flow}}^{-1}+p^{-1}).
\end{equation}
The constant $C$ is independent of the chosen window provided
$p^2\ge W_{\mathrm{flow}}$ and the window exceeds its fixed lower
bound. A kernel vector is reconstructed uniquely from $v=\PP u$,
and its complementary part satisfies
\begin{equation}\label{inf:flow-Q-reconstruction}
 \nrm{\PP^\perp u}\le C\tau W_{\mathrm{flow}}^{-1}
                                      \nrm{\mathsf W v}.
\end{equation}
\end{lemma}
\begin{proof}
Use the positive form normalization $K_0$ and abbreviate
$Q=\PP^\perp$, $\mathcal B_\varsigma=H_\varsigma-H_0$, and
$\mathfrak C_\varsigma=K_0^{-1/2}\mathcal B_{\varsigma}K_0^{-1/2}$.
Write $K_{0,\PP}$ and $K_{0,Q}$ for the restrictions of $K_0$
to $\PP\Hc_{\mathrm{flow}}$ and $Q\Hc_{\mathrm{flow}}$.
The inverse $\mathfrak U_0$ of the normalized unperturbed complement is
radially diagonal with entries
\[
 \frac{z_{j,k}^2+R^2}{z_{j,k}^2-R^2-\lambda},
 \qquad |z_{j,k}^2-R^2|>W_{\mathrm{flow}}.
\]
Their absolute supremum is at most $Cp^2/W_{\mathrm{flow}}$.
This includes the whole infinite tail, where the ratio tends to
one; $p^2\ge W_{\mathrm{flow}}$ covers that limit.
Lemma~\ref{inf:flow-relative} gives
$\nrm{\mathfrak C_\varsigma\mathsf W^{-1}}+
\nrm{\mathfrak C'_\varsigma\mathsf W^{-1}}\le C\tau/p^2$.
Choose the fixed window large enough that
$\nrm{\mathfrak U_0Q\mathfrak C_{\varsigma}Q}\le1/2$ uniformly for $\tau\le20$.
The normalized inverse
\[
 \mathfrak U=(I+\mathfrak U_0Q\mathfrak C_{\varsigma}Q)^{-1}\mathfrak U_0
\]
has norm at most $Cp^2/W_{\mathrm{flow}}$, and
$\mathfrak U'_\varsigma=-\mathfrak UQ\mathfrak C'_{\varsigma}Q\mathfrak U$ has norm at most
$Cp^2\tau/W_{\mathrm{flow}}^2$.
The actual form inverse is $K_{0,Q}^{-1/2}\mathfrak UK_{0,Q}^{-1/2}$.

Eliminate $Q$ using this inverse. The remainder in
\eqref{inf:flow-schur} is
\[
 \mathcal R_\varsigma=\PP\mathcal R_{2,\varsigma}\PP
 -\PP\mathcal B_{\varsigma}QK_{0,Q}^{-1/2}\mathfrak UK_{0,Q}^{-1/2}
                                      Q\mathcal B_\varsigma\PP.
\]
In the second term the two exterior $K_{0,\PP}^{1/2}$ factors cost
$Cp$ each; each normalized perturbation with its exterior inverse
weight costs $C\tau/p^2$. The weighted norm is therefore
$C\theta^2/W_{\mathrm{flow}}$.
Its derivative either replaces one perturbation by $\mathfrak C'_\varsigma$ or replaces
$\mathfrak U$ by $\mathfrak U'_\varsigma$. The added bound in the latter case is
$C\tau^3/W_{\mathrm{flow}}^2$, absorbed in the same bound.
The first remainder term costs $C\theta^2/p$ by
\eqref{inf:flow-quadratic-relative}. This proves
\eqref{inf:flow-schur-remainder} with no dependence of $C$ on the
window size.

The eliminated component is
$u_Q=-K_{0,Q}^{-1/2}\mathfrak UQ\mathfrak C_\varsigma\PP K_{0,\PP}^{1/2}v$.
Since $K_0\ge cp^2I$, the factors cost respectively
$C/p$, $Cp^2/W_{\mathrm{flow}}$, $C\tau/p^2$ after extracting
$\mathsf Wv$, and $Cp$. This proves
\eqref{inf:flow-Q-reconstruction}. The reconstruction belongs to
the common form domain. The eliminated equation is exactly zero
there, so the reduced kernel equation is equivalent to the full
weak eigenvalue equation. A weak solution of that equation is in
the operator domain by the definition of the associated closed-form
operator, completing the equivalence of the reduced and full
eigenvalue equations.
\end{proof}

\subsection{The seed and the second low state}

Assume now that $R$ and $\theta$ satisfy
\eqref{inf:detuning-definition} with $J_p(R)=0$.

\begin{lemma}[Low-block elimination]\label{inf:flow-low-block}
For every fixed $W_{\mathrm{flow}}$ above the absolute lower bound in
Lemma~\ref{inf:flow-complement} and all sufficiently large
$p$, the only angular indices below eight in the flow window are
$j=2,5$, each with one radial state. Their ball offsets satisfy
\begin{equation}\label{inf:low-offsets}
 d_2=-\frac{\theta}{2p}+O(\theta^2p^{-3}),\qquad
 d_5=-72+O(p^{-1}).
\end{equation}
Let $L$ be their two-dimensional span and $\PP_H=\PP-L$.
Write $F_L,\mathsf W_L,u_L$ for the $L$ blocks of $F_\PP,\mathsf W,u$,
and $\mathcal R_{HL},\mathcal R_{LH}$ for the corresponding off-diagonal
blocks of $\mathcal R_\varsigma$ between $L$ and $\PP_H$.
Write $D_H,Z_H,\mathsf W_H$ for the restrictions of
$D_\PP,Z_\PP,\mathsf W$ to this high block, and
$T_H^{\mathrm{ang}}$ for its angular principal submatrix.
For $|\lambda|\le\omega_p=\tau/(16p)$, the $L$ block of
$F_\PP(\varsigma,\lambda)$ is invertible. Eliminating it gives
\begin{equation}\label{inf:high-schur}
 F_H(\varsigma,\lambda)=D_H-\lambda+B_{1,HH}
                                      +\widehat{\mathcal R}_\varsigma(\lambda),
\end{equation}
with
\begin{equation}\label{inf:high-schur-remainder}
 \nrm{\mathsf W_H^{-1}(\partial_\varsigma^k\widehat{\mathcal R}_\varsigma)
                         \mathsf W_H^{-1}}
 \le C\theta^2(W_{\mathrm{flow}}^{-1}+p^{-1}),\quad k=0,1.
\end{equation}
A nonzero full eigenvector in this band has $v=\PP_Hu\ne0$ and
$\nrm{u}\le C\nrm v$, uniformly also as $\tau\downarrow0$.
\end{lemma}
\begin{proof}
The trace ratio and \eqref{inf:rho-equation} give
$R=2p+3+O(p^{-1})$. Normalize the seed radial function by
$f'(R)=1$, $f(R)=\chi=\theta/(4Rp)$.
The collar equation has bounded coefficients and one bounded
derivative for fixed angular index. On $|y|\le2|\chi|$ its integral
equation yields $f'(R+y)=1+O(|\chi|)$. The endpoint values have
opposite signs, and the derivative has positive sign.
Its unique zero in the collar is
$R-\chi+O(\chi^2)$. Squaring gives
$d_2=-2R\chi+O(R\chi^2)=-\theta/(2p)+O(\theta^2p^{-3})$,
uniformly for arbitrarily small nonzero $\theta$.

For $j=5$ use \eqref{inf:resonant-trace} with $c_5=18$.
The same collar argument places its zero at
$R-18/p+O(p^{-2})$, giving the second offset in
\eqref{inf:low-offsets}. For $0\le j<8$ other than $2,5$, the
finite recurrence \eqref{inf:finite-q-recurrence} has nonzero
limiting value at $k=2j-1$. Its derivative trace stays bounded.
The collar equation excludes a zero in a fixed positive-width
collar of $R$. Such an index is therefore outside every fixed
energy window for large $p$. The radial spacing excludes a second
state in each remaining index.

The first-order operator has no $L$--$\PP_H$ coupling or coupling
between the two $L$ states. Their angular distances are at least
three, whereas both matrices in
\eqref{inf:transport-identity} have angular bandwidth two.
The seed diagonal is
\[
 \sigma_2=d_2+\frac{\theta z_2^2}{pR^2}
          =\frac{\theta}{2p}+O(\theta^2p^{-3}),
\]
because $(T_\varsigma)_{22}=1/p$ exactly. The other diagonal tends to $-72$
and its perturbation is $O(\tau/p)$.
For $|\lambda|\le\omega_p$, these two uncorrected diagonal entries
have modulus at least $\tau/(4p)$ and $36$, respectively.
The remainder block has norm at most
$C\theta^2(W_{\mathrm{flow}}^{-1}+p^{-1})p^{-2}$ because
$w_2,w_5\le6/p$. A two-dimensional Neumann series gives
\[
 \nrm{F_L^{-1}}+\nrm{\partial_{\varsigma}F_L^{-1}}\le Cp/\tau.
\]
For the derivative use
$(F_L^{-1})'=-F_L^{-1}F'_LF_L^{-1}$ and
$\nrm{F'_L}\le C\tau/p$, which follows from the diagonal structure
and the angular derivative bound. In particular,
\begin{equation}\label{inf:weighted-low-inverse}
 \nrm{\mathsf W_LF_L^{-1}\mathsf W_L}+
 \nrm{\mathsf W_L(\partial_{\varsigma}F_L^{-1})\mathsf W_L}
 \le C/(p\tau).
\end{equation}
The second elimination has remainder
\[
 \widehat{\mathcal R}
 =\mathcal R_{HH}-\mathcal R_{HL}F_L^{-1}\mathcal R_{LH}.
\]
Put $a_*=C\theta^2(W_{\mathrm{flow}}^{-1}+p^{-1})$.
The weighted new term and its derivative are at most
$Ca_*^2/(p\tau)$ by \eqref{inf:weighted-low-inverse}.
For $\tau\le20$ this is absorbed in $a_*$ after increasing the
$p$ threshold, proving \eqref{inf:high-schur-remainder}.
Reconstruction gives
\[
 \nrm{u_L}\le C\tau(W_{\mathrm{flow}}^{-1}+p^{-1})
                           \nrm{\mathsf W_Hv}.
\]
Combining this with \eqref{inf:flow-Q-reconstruction} proves
$\nrm u\le C\nrm v$ and the nonvanishing assertion.
The inverse factors $1/\tau$ are paired with the quadratic
remainders, giving uniform reconstruction for $0<\tau\le20$.
\end{proof}

\begin{lemma}[Ordered eigenvalue derivatives]\label{inf:ordered-eigenvalues}
For the common $C^1$ closed-form family above, the ordered
eigenvalues are locally Lipschitz. At an interior parameter
$\varsigma_0\in(1/64,1/16)$, suppose an eigenvalue $\lambda$ has multiplicity
$m$ and occupies ordered positions $i,\ldots,i+m-1$.
Let $\mu_1\le\cdots\le\mu_m$ be the eigenvalues of the derivative
form restricted to its Hilbert eigenspace. The one-sided
derivatives exist and satisfy
\begin{equation}\label{inf:one-sided-eigenvalues}
 \lambda'_{i+r-1,+}(\varsigma_0)=\mu_r,\qquad
 \lambda'_{i+r-1,-}(\varsigma_0)=\mu_{m-r+1},\qquad 1\le r\le m.
\end{equation}
Thus every one-sided derivative is an eigenvalue of the
compressed derivative form. An ordered eigenvalue need not
have a two-sided derivative at a multiple crossing.
\end{lemma}
\begin{proof}
A fixed positive shift has inverse that is norm $C^1$, by the
inverse identity on the common form space. Its compact resolvent
admits a contour surrounding the isolated finite cluster in
question. The contour integral gives a $C^1$ finite-rank
projection, also with values in the form domain. Local unitary
transport of this projection identifies its range with a fixed
finite-dimensional space. Restricting the resolvent there and
inverting that finite matrix gives a $C^1$ self-adjoint matrix
for the original spectral cluster; see also
\cite[Chapter~II, Theorem~5.4 and Theorem~6.8]{inf:kato}.
At $\varsigma_0$ it is $\lambda I$. Its derivative is the
restricted derivative form: differentiation of its unitary
transport contributes a commutator with $\lambda I$, which is
zero. The matrix expansion is therefore $\lambda I+hB+o(|h|)$,
with $B$ having eigenvalues $\mu_r$.
The finite-dimensional min--max principle bounds changes of
ordered eigenvalues by the operator norm of the remainder.
For $h>0$ the order of $h\mu_r$ is unchanged, and for $h<0$ it
is reversed. Division by $h$ gives
\eqref{inf:one-sided-eigenvalues}.
The same estimate for two nearby parameters gives local
Lipschitz continuity. Compactness of the parameter interval
gives absolute continuity for each fixed ordered curve, hence
a two-sided derivative almost everywhere.
\end{proof}

\begin{theorem}[Local spectral flow]\label{inf:flow-theorem}
For all sufficiently large $p\in\mathcal P_{\mathrm{bd}}$, let $R,\theta$
be as in Lemma~\ref{inf:detuning}. There are absolute constants
$c_{\mathrm{flow}}>0$ and $P_{\mathrm{flow}}<\infty$ such that
every ordered eigenvalue offset $\lambda_i(\varsigma)$ of $H_\varsigma-R^2$
satisfies
\begin{equation}\label{inf:flow-conclusion}
 \operatorname{sign}(\theta)\lambda_i'(\varsigma)
 \ge c_{\mathrm{flow}}\tau/p
 \quad\hbox{for almost every }\varsigma\hbox{ with }
 |\lambda_i(\varsigma)|<\tau/(16p),
\end{equation}
when $p\ge P_{\mathrm{flow}}$. The constants are uniform in
$0<\tau\le20$.
\end{theorem}
\begin{proof}
Fix an eigenvalue $\lambda$ in the band and a vector $u$ in its
entire eigenspace. Put $v=\PP_Hu$. The two exact eliminations show
that $v\ne0$. With $v,\lambda$ held fixed, differentiation of the
reconstructed quadratic form gives
\begin{equation}\label{inf:feshbach-derivative}
 \langle u,H'_{\varsigma}u\rangle
 =\langle v,\partial_{\varsigma}F_H(\varsigma,\lambda)v\rangle.
\end{equation}
Indeed, terms differentiating the reconstruction lie in
$\PP^\perp+L$, and their pairing with the operator vanishes by the
eliminated equations. The bounded normalized inverses above
make the reconstruction differentiable in the common form norm.

The transport identity gives
\[
 \partial_{\varsigma}F_H
 =\theta Z_H(T_H^{\mathrm{ang}})'Z_H
          -[D_H,\mathcal S'_{HH}]+\partial_\varsigma\widehat{\mathcal R}.
\]
The first term, after multiplication by $\operatorname{sign}\theta$,
is at least $c_1\tau\nrm{\mathsf W_Hv}^2$ in quadratic form.
Here Lemma~\ref{inf:angular-positive} applies to all the retained
indices, and $z_i/R\to1$ uniformly in the fixed window; $c_1>0$
is absolute.
Estimate the transport term on the actual kernel vector, using
\eqref{inf:high-schur}:
$D_Hv=\lambda v-(B_{1,HH}+\widehat{\mathcal R})v$.
Since $\mathcal S'_{HH}$ is skew-adjoint, its real pairing with $\lambda v$
vanishes. Thus
\begin{align*}
 |\langle v,[D_H,\mathcal S'_{HH}]v\rangle|
 &\le2\nrm{(B_{1,HH}+\widehat{\mathcal R})v}
                                  \nrm{\mathcal S'_{HH}v}\\
 &\le C\theta^2p^{-1}\nrm{\mathsf W_Hv}^2.
\end{align*}
The final remainder is at most
$C\theta^2(W_{\mathrm{flow}}^{-1}+p^{-1})
\nrm{\mathsf W_Hv}^2$.
Choose one constant $C_{\mathrm{flow,err}}$ dominating the
transport-error constant and the constants in
\eqref{inf:flow-schur-remainder} and
\eqref{inf:high-schur-remainder}. It is independent of
$W_{\mathrm{flow}}$, under their stated lower bound on that
window. The combined error is at most
$C_{\mathrm{flow,err}}\theta^2(W_{\mathrm{flow}}^{-1}+2p^{-1})
\nrm{\mathsf W_Hv}^2$.
Choose $W_{\mathrm{flow}}$ so large that
$20C_{\mathrm{flow,err}}/W_{\mathrm{flow}}<c_1/4$ and the
complementary-inverse conditions hold. Next choose $p$ large
enough that $40C_{\mathrm{flow,err}}/p<c_1/4$.
The other window-dependent requirements are
$p^2\ge W_{\mathrm{flow}}$, $2W_{\mathrm{flow}}/R<1/10$,
and the fixed-index asymptotic threshold in
Lemma~\ref{inf:flow-low-block}. Also impose the near-crossing
low-degree threshold in Proposition~\ref{inf:low-exclusion}. Equations
\eqref{inf:feshbach-derivative} and these bounds give
\begin{equation}\label{inf:weighted-flow-margin}
 \operatorname{sign}(\theta)\langle u,H'_{\varsigma}u\rangle
 \ge(c_1/2)\tau\nrm{\mathsf W_Hv}^2.
\end{equation}

To turn the weighted bound into a $1/p$ margin, put
$J_p=\lfloor\sqrt p/16\rfloor$.
For large $p$ every index $j<J_p$ other than the removed seed is in
the exclusion range of Proposition~\ref{inf:low-exclusion}.
Its ball offset has modulus greater than $32$.
On those rows of \eqref{inf:high-schur}, invert $D_H-\lambda$,
whose inverse norm is at most $1/31$, to get
\[
 \nrm{v_{j<J_p}}\le C\tau\nrm{\mathsf W_Hv}.
\]
On the other indices $w_j\ge c/\sqrt p$, giving
$\nrm{v_{j\ge J_p}}\le C\sqrt p\nrm{\mathsf W_Hv}$.
Lemma~\ref{inf:flow-low-block} therefore yields
$\nrm{\mathsf W_Hv}^2\ge cp^{-1}\nrm u^2$.
The further resonant low states $j=8,11,\ldots$ in a large fixed
window are included in this argument: they remain in $\PP_H$,
are covered by its positive angular variation, and are excluded
from the near-zero mass by the last row estimate.
Substitution into \eqref{inf:weighted-flow-margin} proves the
claimed derivative lower bound on the entire eigenspace.

By Lemma~\ref{inf:ordered-eigenvalues}, every one-sided derivative
of an ordered eigenvalue in the band is an eigenvalue of this
compressed derivative form. Its signed lower bound therefore
applies at every differentiability point, and almost everywhere
on each absolutely continuous curve. This proves
\eqref{inf:flow-conclusion}, including multiple crossings.
\end{proof}
